\documentclass[aspectratio=169,11pt]{beamer}

\usetheme{Madrid}
\usecolortheme{default}
\usefonttheme{professionalfonts}
\setbeamertemplate{navigation symbols}{}
\setbeamertemplate{footline}[frame number]

\usepackage{amsmath, amssymb, booktabs}

\newcommand{\E}{\mathbb{E}}
\newcommand{\1}{\mathbf{1}}
\newcommand{\Prb}{\mathbb{P}}

\title{Spatial Structural Modeling}
\subtitle{Empirical Methods --- Week 17}
\author{Teaching deck draft}
\date{}

\begin{document}

\begin{frame}
  \titlepage
\end{frame}

\begin{frame}{Week 17 question}
\begin{itemize}
    \item When do spatial research questions require equilibrium structure?
    \item Lecture 15 built spatial objects: places, distances, exposures, and access.
    \item Lecture 16 asked what identifies causal effects when space creates spillovers, sorting, or correlated shocks.
    \item Lecture 17 asks when workers, firms, rents, amenities, commuting, migration, and welfare must be modeled jointly.
\end{itemize}

\medskip
Structure is useful when the policy object is an equilibrium outcome, not just a local treatment effect.
\end{frame}

\begin{frame}{When reduced form is not enough}
\small
\begin{tabular}{p{0.29\linewidth} p{0.30\linewidth} p{0.30\linewidth}}
\toprule
Reduced-form fact & Missing equilibrium object & Why it matters \\
\midrule
Local wages rise & rents, amenities, commuting costs & nominal wages are not welfare \\
Employment changes in treated place & migration and commuting flows & local shocks do not stay local \\
Housing reform affects prices & housing supply and sorting & gains may accrue to workers or landlords \\
Transit lowers travel time & job access and firm response & residence and workplace incidence differ \\
Place subsidy raises jobs & firm location and spillovers & nearby places may lose or gain \\
\bottomrule
\end{tabular}

\medskip
\[
\text{Need structure when } \mathcal{E}(d)
=
\{w,r,F,L,H,U\}(d)
\text{ is the object.}
\]
\end{frame}

\begin{frame}{Core equilibrium map}
\small
\[
V_{ioj}
=
w_j-r_i-c_{ij}^{commute}-m_{oi}+a_i+\varepsilon_{ioj}
\]

\[
P_{ioj}
=
\frac{\exp(\bar V_{ioj}/\sigma)}
{\sum_{i'}\sum_{j'}\exp(\bar V_{oi'j'}/\sigma)}
\]

\[
\mathcal{E}(\theta;d)
=
\left\{
P_{ioj},\pi_{ij}^{commute},w_j,r_i,L_j,H_i,U_o
\right\}_{o,i,j}
\]

\begin{itemize}
    \item Workers choose locations and commutes given wages, rents, amenities, and frictions.
    \item Firms demand labor where productivity and market access are high.
    \item Housing supply and congestion determine rent incidence.
    \item Equilibrium closes through prices, quantities, and bilateral flows.
\end{itemize}
\end{frame}

\begin{frame}{Core objects to keep separate}
\small
\begin{tabular}{p{0.24\linewidth} p{0.36\linewidth} p{0.28\linewidth}}
\toprule
Object & Typical evidence & Main risk \\
\midrule
Moving costs & migration flows, mover histories & confuse amenities with frictions \\
Commuting costs & residence-workplace flows, travel time & treat distance as generalized cost \\
Local productivity & wages, output, firm data & absorb amenities or sorting \\
Amenities & residual utility, housing demand & residual becomes everything \\
Rents and housing supply & prices, quantities, regulation, land & miss landlord incidence \\
Labor demand and firm location & employment, establishments, shocks & hold firms fixed when policy moves them \\
\bottomrule
\end{tabular}

\medskip
The model is credible only if these objects are not silently bundled together.
\end{frame}

\begin{frame}{Static versus dynamic spatial structure}
\small
\begin{tabular}{p{0.30\linewidth} p{0.30\linewidth} p{0.30\linewidth}}
\toprule
Question & Static model & Dynamic model \\
\midrule
Main use & long-run incidence and comparative statics & transition paths and policy timing \\
Adjustment & workers, firms, rents, flows clear after shock & migration, housing, tenure, and investment adjust over time \\
Data pressure & cross-section, long differences, flows & panels, histories, timing, expectations \\
Best when & long-run allocation is the object & adjustment costs are the object \\
Main caveat & hides transition costs & heavier assumptions and computation \\
\bottomrule
\end{tabular}

\medskip
The static-dynamic choice should follow the estimand, not the desire for a larger model.
\end{frame}

\begin{frame}{Gravity-style building blocks}
\small
\[
F_{ij}
=
\frac{O_iD_j\phi_{ij}^{-\kappa}}
{\sum_{\ell}D_{\ell}\phi_{i\ell}^{-\kappa}}
\qquad
\log F_{ij}
=
\alpha_i+\delta_j-\kappa\log \phi_{ij}+e_{ij}
\]

\begin{itemize}
    \item Gravity disciplines bilateral trade, commuting, migration, applications, or trips.
    \item Observed flows help recover frictions and flow elasticities.
    \item The same algebra can sit inside many spatial models.
    \item Gravity is a building block, not the whole model, when the policy question is incidence or welfare.
\end{itemize}
\end{frame}

\begin{frame}{Identification, calibration, and fit}
\small
\begin{tabular}{p{0.24\linewidth} p{0.35\linewidth} p{0.29\linewidth}}
\toprule
Model object & Empirical discipline & Diagnostic \\
\midrule
Commuting friction & bilateral residence-workplace flows & gravity residuals, untargeted flows \\
Migration cost & mover flows and histories & response to shocks, origin persistence \\
Productivity & wages, output, firm data & separate from amenities and composition \\
Housing supply & rents, quantities, regulation & sensitivity to elasticity \\
Amenities & residual utility or external measures & validate with non-targeted outcomes \\
Counterfactual & policy shock and fixed point & distance from observed support \\
\bottomrule
\end{tabular}

\medskip
Report what is estimated, what is calibrated, what is inverted, and what is held fixed.
\end{frame}

\begin{frame}{Counterfactual interpretation}
\small
\[
\Delta W
=
\sum_g \omega_g
\left[
U_g(\mathcal{E}(\theta;d_1))
-
U_g(\mathcal{E}(\theta;d_0))
\right]
\]

\begin{itemize}
    \item Welfare requires a utility object, not only wages or employment.
    \item Incidence can differ for residents, in-commuters, movers, firms, landlords, and taxpayers.
    \item Nearby policy changes are more credible than large reforms far outside support.
    \item Sensitivity should target the margins that move the headline result: frictions, housing supply, amenities, and firm response.
\end{itemize}
\end{frame}

\begin{frame}{Spatial methods block summary}
\small
\begin{tabular}{p{0.20\linewidth} p{0.33\linewidth} p{0.35\linewidth}}
\toprule
Lecture & Skill & Research payoff \\
\midrule
15 & curate spatial objects & places, distances, crosswalks, access, exposure \\
16 & identify effects with space & spillovers, sorting, borders, shift-share, inference \\
17 & model equilibrium propagation & wages, rents, flows, firms, housing, welfare \\
\bottomrule
\end{tabular}

\medskip
Where the literature is now:
\begin{itemize}
    \item richer bilateral flow data and quantitative spatial models;
    \item more attention to equilibrium properties, multiplicity, and computation;
    \item dynamic urban models with transition paths, housing frictions, and forward-looking households;
    \item frontier questions on remote work, climate adaptation, monopsony in space, and heterogeneous workers.
\end{itemize}
\end{frame}

\begin{frame}{Research Lab design}
\begin{columns}[T,onlytextwidth]
\begin{column}{0.32\linewidth}
\textbf{Reproduce}
\begin{itemize}
    \item synthetic locations
    \item commuting gravity
    \item commuting openness
    \item local productivity shock
\end{itemize}
\end{column}
\begin{column}{0.32\linewidth}
\textbf{Diagnose}
\begin{itemize}
    \item targeted vs untargeted fit
    \item friction sensitivity
    \item housing supply sensitivity
    \item reduced-form stopping point
\end{itemize}
\end{column}
\begin{column}{0.32\linewidth}
\textbf{Transfer}
\begin{itemize}
    \item dynamic adjustment path
    \item migration friction
    \item gradual housing response
    \item static versus dynamic memo
\end{itemize}
\end{column}
\end{columns}

\medskip
\small
Primary anchor: Monte--Redding--Rossi-Hansberg. Dynamic extension: Greaney--Parkhomenko--Van Nieuwerburgh.
\end{frame}

\begin{frame}{Takeaways}
\begin{itemize}
    \item Spatial structural models are justified by the object, not by style.
    \item The core object is an equilibrium mapping from shocks to prices, flows, quantities, and welfare.
    \item Gravity-style blocks are valuable flow discipline inside that mapping.
    \item Identification, calibration, fit, and counterfactual distance should be reported separately.
    \item The spatial block now gives three tools: build spatial objects, identify spatial effects, and model equilibrium propagation.
\end{itemize}
\end{frame}

\end{document}
